\answer{ex:06:1}{$e=(1-e^{-T_a/\tau_e})e^{-d/\tau_e}$. (a)~$0.110$ so với $4.5\cdot10^{-5}$, gấp $\approx2.5\cdot10^3$ lần. (b)~$\tau_e^*=T_a/\ln(1+T_a/d)=0.59\,\text{s}$.}
\answer{ex:06:2}{Tốc độ $\nu_x\nu_y(A_+\tau_+-A_-\tau_-)=0.8A_+$ mỗi giây, $\approx2900A_+$ mỗi giờ; $\tau_-=\tau_+/0.6=33\ms$.}
\answer{ex:06:3}{$\mathbf e_1$ khi $\mu^2+s_1^2>s_2^2$; ở đây $1.01>0.25$, và nơ-ron học được $\mathbf e_1$, dù độ phân tán theo $\mathbf e_2$ lớn hơn $25$ lần: quy tắc tìm vectơ chính của ma trận tương quan chứ không phải ma trận hiệp phương sai.}
\answer{ex:06:4}{$V=Re^{-(t_c+L-t)/\tau_\gamma}$ giữa lúc đèn sáng và lúc có nước ép, nên $\dot V=V/\tau_\gamma$; đỉnh vọt có diện tích $Re^{-L/\tau_\gamma}$: $0.61R$ và $0.78R$.}
\answer{ex:06:5}{Tỷ số tín hiệu trên nhiễu là $1/(N+1)$, số lượt thử $\propto N+1$: $5\cdot10^5$ lượt thử $\approx1$~tháng và $5\cdot10^{11}$ lượt thử $\approx8\cdot10^4$~năm.}
\answer{ex:06:6}{(a)~$k_2=1/\tau_d$, $k_1=1/\tau_d-1/\tau_\gamma$. (b)~Sai số giả $-(V/\tau_\gamma)\,\varepsilon/(1+\varepsilon)$, $\varepsilon=\tau_d/\tau_\gamma$, suy ra $\tau_d\le0.105\,\text{s}$.}
\answer{ex:06:7}{$0.46$, $0.90$, $0.99$, $0.95$ và $0.70$ khi $d=3\,\text{s}$. Các điểm nằm trên cùng một đường cong theo phần vết còn lại $F=(1-e^{-T_{\text{cue}}/\tau_e})e^{-d/\tau_e}$: học được khi $F\gtrsim0.1$, chọn ngẫu nhiên khi $F<10^{-2}$.}
\answer{ex:06:8}{$\eta^*=1/\bigl(2\sigma^2(N+2)\bigr)$, sau $N+2$ lượt thử; khi chỉ $m$ trong $N$ nơ-ron dao động --- sau $N(m+2)/m\ge N+2$: tính thưa không giúp ích.}
